\par\Needspace{9\baselineskip}\bigskip
\pdfbookmark[3]{Exposition, [24]}{pallas-II-article24}
\begin{center}[24].\end{center}
\begingroup
\makeatletter
\def\@makefnmark{\hbox{\normalfont\@thefnmark}}
\makeatother
\noindent\hspace*{\parindent}If the ratio of \(n\text{ and }n'\) is rational, the value of \(n+\delta n\) will contain a part proportional to the time, or a secular variation, and the mean motion \(\int n\mathrm{d}t\) will contain a term proportional to the square of the time. We can leave it at that for the present needs of astronomy, but we shall see that considering terms that depend on higher powers of the masses restores a \textit{periodic form} to this apparently secular equation. This is all the more remarkable because the theory of the motion of Pallas will in fact furnish us with this case\HistoricalNoteEight{See pp. 557---559.] }.
\par\endgroup
